% GENERATED FILE. Edit canonical lessons, then run npm run generate:books.
% canonical-source-hash: fnv1a64:a7a975a944e3a698
% canonical-lessons: GU-C309-sabha, GU-C309-sanman, GU-C309-inam, GU-C309-birdavvu, GU-C309-garva

\chapter{An Award for Courage}
\label{ch:gu-an-award-for-courage}
\hlchaptermodality{309}

\begin{chapteropening}
\textbf{By the end of this chapter:} \emph{I can say who was honoured and why.}

Complete the last lesson of chapter 309: I can say who was honoured and why.
\end{chapteropening}

\section[\texorpdfstring{\gu{સભા} (sabhā)}{સભા (sabhā)}]{\gu{સભા} (sabhā) — a gathering; an assembly}

\label{lesson:GU-C309-sabha}



\begin{quote}
Before the new one: say the Gujarati for patience, then the Gujarati for to plead.
\end{quote}

\subsection*{You'll want to know: \gu{સભા}}
\textbf{\gu{સભા}} — \emph{sabhā} — \textquotedblleft{}a gathering; an assembly\textquotedblright{}.

\begin{quote}
\gu{સાંજે} \gu{ગામમાં} \gu{મોટી} \gu{સભા} \gu{થઈ}.
\end{quote}

\begin{quote}
\emph{sā̃je gāmmā̃ moṭī sabhā thaī.}
\end{quote}

\begin{quote}
In the evening there was a big gathering in the village.
\end{quote}

A feminine noun.

\begin{cousinweb}
is a tatsama, Sanskrit \emph{sabhā}, \textquotedblleft{}an assembly, a meeting place\textquotedblright{} (Monier-Williams). A \textbf{\gu{સભા}} is people gathered to hear and to decide; \textbf{\gu{બેઠક}} is a smaller meeting.
\end{cousinweb}

\begin{grammarlens}[title={What you've built}]
One more word for telling what happened.
\end{grammarlens}

\subsection*{Your turn}
\begin{itemize}
\raggedright
  \item \emph{Say it:} \emph{sabhā}
  \item \emph{Say it:} \emph{sabhā}, once more
  \item \emph{Say it:} \emph{sā̃je gāmmā̃ moṭī sabhā thaī.} — In the evening there was a big gathering in the village.
  \item \emph{Recall:} say the Gujarati for patience, then the Gujarati for to plead, then say \emph{sabhā} again
\end{itemize}

\begin{tcolorbox}[breakable,colback=teal!4,colframe=teal!35!black,title={Before you move on}]
What does \gu{સભા} mean? (\textquotedblleft{}A gathering; an assembly\textquotedblright{}.) Say it once more.
\end{tcolorbox}
\section[\texorpdfstring{\gu{સન્માન} (sanmān)}{સન્માન (sanmān)}]{\gu{સન્માન} (sanmān) — an honour}

\label{lesson:GU-C309-sanman}



\begin{quote}
Before the new one: say the Gujarati for to plead, then the Gujarati for an assembly.
\end{quote}

\subsection*{You'll want to know: \gu{સન્માન}}
\textbf{\gu{સન્માન}} — \emph{sanmān} — \textquotedblleft{}an honour\textquotedblright{}.

\begin{quote}
\gu{સભામાં} \gu{અરુણનું} \gu{સન્માન} \gu{થયું}.
\end{quote}

\begin{quote}
\emph{sabhāmā̃ Aruṇnũ sanmān thayũ.}
\end{quote}

\begin{quote}
Arun was honoured at the gathering.
\end{quote}

A neuter noun.

\begin{cousinweb}
\textbf{\gu{સન્માન}} (\emph{sanmān}) is a tatsama, Sanskrit \emph{sammāna}, \textquotedblleft{}honour, respect, homage\textquotedblright{} (Monier-Williams). Someone is honoured with \textbf{\gu{થવું}}: \textbf{\gu{તેનું} \gu{સન્માન} \gu{થયું}}.
\end{cousinweb}

\begin{grammarlens}[title={What you've built}]
One more word for telling what happened.
\end{grammarlens}

\subsection*{Your turn}
\begin{itemize}
\raggedright
  \item \emph{Say it:} \emph{sanmān}
  \item \emph{Say it:} \emph{sanmān}, once more
  \item \emph{Say it:} \emph{sabhāmā̃ Aruṇnũ sanmān thayũ.} — Arun was honoured at the gathering.
  \item \emph{Recall:} say the Gujarati for to plead, then the Gujarati for an assembly, then say \emph{sanmān} again
\end{itemize}

\begin{tcolorbox}[breakable,colback=teal!4,colframe=teal!35!black,title={Before you move on}]
What does \gu{સન્માન} mean? (\textquotedblleft{}An honour\textquotedblright{}.) Say it once more.
\end{tcolorbox}
\section[\texorpdfstring{\gu{ઇનામ} (inām)}{ઇનામ (inām)}]{\gu{ઇનામ} (inām) — a prize; a reward}

\label{lesson:GU-C309-inam}



\begin{quote}
Before the new one: say the Gujarati for an assembly, then the Gujarati for an honour.
\end{quote}

\subsection*{You'll want to know: \gu{ઇનામ}}
\textbf{\gu{ઇનામ}} — \emph{inām} — \textquotedblleft{}a prize; a reward\textquotedblright{}.

\begin{quote}
\gu{ગામના} \gu{લોકોએ} \gu{અરુણને} \gu{ઇનામ} \gu{આપ્યું}.
\end{quote}

\begin{quote}
\emph{gāmnā lokoe Aruṇne inām āpyũ.}
\end{quote}

\begin{quote}
The villagers gave Arun a reward.
\end{quote}

A neuter noun.

\begin{cousinweb}
\textbf{\gu{ઇનામ}} (\emph{inām}) is usually traced, through Persian, to Arabic \emph{inʿām}, \textquotedblleft{}something bestowed\textquotedblright{}. It is neuter: \textbf{\gu{ઇનામ} \gu{આપ્યું}}, gave a prize; \textbf{\gu{ઇનામ} \gu{જીત્યું}}, won one.
\end{cousinweb}

\begin{grammarlens}[title={What you've built}]
One more word for telling what happened.
\end{grammarlens}

\subsection*{Your turn}
\begin{itemize}
\raggedright
  \item \emph{Say it:} \emph{inām}
  \item \emph{Say it:} \emph{inām}, once more
  \item \emph{Say it:} \emph{gāmnā lokoe Aruṇne inām āpyũ.} — The villagers gave Arun a reward.
  \item \emph{Recall:} say the Gujarati for an assembly, then the Gujarati for an honour, then say \emph{inām} again
\end{itemize}

\begin{tcolorbox}[breakable,colback=teal!4,colframe=teal!35!black,title={Before you move on}]
What does \gu{ઇનામ} mean? (\textquotedblleft{}A prize; a reward\textquotedblright{}.) Say it once more.
\end{tcolorbox}
\section[\texorpdfstring{\gu{બિરદાવવું} (birdāvvũ)}{બિરદાવવું (birdāvvũ)}]{\gu{બિરદાવવું} (birdāvvũ) — to applaud; to praise publicly}

\label{lesson:GU-C309-birdavvu}



\begin{quote}
Before the new one: say the Gujarati for an honour, then the Gujarati for a prize.
\end{quote}

\subsection*{You'll want to know: \gu{બિરદાવવું}}
\textbf{\gu{બિરદાવવું}} — \emph{birdāvvũ} — \textquotedblleft{}to applaud; to praise publicly\textquotedblright{}.

\begin{quote}
\gu{બધાએ} \gu{અરુણના} \gu{સાહસને} \gu{બિરદાવ્યું}.
\end{quote}

\begin{quote}
\emph{badhāe Aruṇnā sāhasne birdāvyũ.}
\end{quote}

\begin{quote}
Everyone applauded Arun's courage.
\end{quote}

Past, with \gu{તેણે} (he or she …): \gu{તેણે} \gu{તેને} \gu{બિરદાવ્યો}, \emph{teṇe tene birdāvyo}.

\begin{cousinweb}
is built on \emph{birad}, praise, from Sanskrit \emph{viruda}, \textquotedblleft{}a laudatory poem\textquotedblright{}, which Monier-Williams also gives as \emph{birada}. \textbf{\gu{વખાણવું}} is to praise; \textbf{\gu{બિરદાવવું}} is to praise in public, before others.
\end{cousinweb}

\begin{grammarlens}[title={What you've built}]
One more word for telling what happened.
\end{grammarlens}

\subsection*{Your turn}
\begin{itemize}
\raggedright
  \item \emph{Say it:} \emph{birdāvvũ}
  \item \emph{Say it:} \emph{birdāvvũ}, once more
  \item \emph{Say it:} \emph{badhāe Aruṇnā sāhasne birdāvyũ.} — Everyone applauded Arun's courage.
  \item \emph{Recall:} say the Gujarati for an honour, then the Gujarati for a prize, then say \emph{birdāvvũ} again
\end{itemize}

\begin{tcolorbox}[breakable,colback=teal!4,colframe=teal!35!black,title={Before you move on}]
What does \gu{બિરદાવવું} mean? (\textquotedblleft{}To applaud; to praise publicly\textquotedblright{}.) Say it once more.
\end{tcolorbox}
\section[\texorpdfstring{\gu{ગર્વ} (garva)}{ગર્વ (garva)}]{\gu{ગર્વ} (garva) — pride}

\label{lesson:GU-C309-garva}



\begin{quote}
Before the new one: say the Gujarati for a prize, then the Gujarati for to applaud.
\end{quote}

\subsection*{You'll want to know: \gu{ગર્વ}}
\textbf{\gu{ગર્વ}} — \emph{garva} — \textquotedblleft{}pride\textquotedblright{}.

\begin{quote}
\gu{માતાને} \gu{દીકરા} \gu{પર} \gu{ખૂબ} \gu{ગર્વ} \gu{હતો}.
\end{quote}

\begin{quote}
\emph{mātāne dīkrā par khūb garva hato.}
\end{quote}

\begin{quote}
The mother was very proud of her son.
\end{quote}

A masculine noun.

\begin{cousinweb}
\textbf{\gu{ગર્વ}} (\emph{garva}) is a tatsama, Sanskrit \emph{garva}, \textquotedblleft{}pride, arrogance\textquotedblright{} (CDIAL 4064). In Gujarati it is often good pride: \textbf{\gu{દીકરા} \gu{પર} \gu{ગર્વ}}, proud of a son.
\end{cousinweb}

\begin{grammarlens}[title={What you've built}]
One more word for telling what happened.
\end{grammarlens}

\subsection*{Your turn}
\begin{itemize}
\raggedright
  \item \emph{Say it:} \emph{garva}
  \item \emph{Say it:} \emph{garva}, once more
  \item \emph{Say it:} \emph{mātāne dīkrā par khūb garva hato.} — The mother was very proud of her son.
  \item \emph{Recall:} say the Gujarati for a prize, then the Gujarati for to applaud, then say \emph{garva} again
\end{itemize}

\begin{tcolorbox}[breakable,colback=teal!4,colframe=teal!35!black,title={Before you move on}]
What does \gu{ગર્વ} mean? (\textquotedblleft{}Pride\textquotedblright{}.) Say it once more.
\end{tcolorbox}
